Por objetivo, tomadas del §\ref{sec:estado-arte} y de los factores del Capítulo~\ref{cap:objetivos}.
No hay aprendizaje supervisado: S y E salen de reglas de \textit{schema matching} y de similitud,
no de un modelo etiquetado.

\begin{itemize}
  \item \textbf{OE1:} \textit{parsing} OpenAPI~3.x; extracción de \textit{paths}, operaciones y
        esquemas; representación intermedia alineada con BIAN (Mainas et al., 2023 como referencia
        de enriquecimiento semántico).
  \item \textbf{OE2:} \textit{schema matching} según taxonomía Shvaiko y Euzenat (2005): nombre,
        ruta, profundidad de esquema, operación HTTP.
  \item \textbf{OE3:} similitud semántica por embeddings y/o ontología (El Bayadh et al., 2015;
        Chandrasekaran y Mago, 2022).
  \item \textbf{OE4:} agregación ponderada $\alpha S + \beta E + \gamma C$ y clasificación por
        umbrales fijos (§1.4).
\end{itemize}

Herramientas previstas: Python para prototipo, validadores OpenAPI, bibliotecas de similitud textual.
La implementación concreta queda fuera del alcance del presente plan (§\ref{sec:alcances}).
